\section{Аналіз базового синтаксису та операторів: порівняльний аспект із C/C++}

\subsection{Загальні синтаксичні конструкції та типи даних}

Оскільки синтаксис мови Java історично базується на мовах C та C++, базові керуючі 
конструкції (цикли \texttt{for}, \texttt{while}, \texttt{do-while} та оператори 
розгалуження \texttt{if-else}, \texttt{switch}) є практично ідентичними. Проте на 
рівні базової архітектури та типізації мови мають суттєві відмінності:

\begin{itemize}
    \item \textbf{Строга типізація та незалежність від платформи:} 
    На відміну від C/C++,  де розмір типу \texttt{int} чи \texttt{long} залежить 
    від архітектури процесора та компілятора, у Java всі примітивні типи мають фіксований 
    розмір у пам'яті на будь-якій платформі (наприклад, \texttt{int} завжди 32 біти, \texttt{long} — 64 біти).
    \item \textbf{Відсутність беззнакових типів:} 
    У Java відсутній модифікатор \texttt{unsigned}. Усі числові примітивні 
    типи (включаючи \texttt{byte} та \texttt{short}) є виключно знаковими.
    \item \textbf{Булевий тип даних:} 
    На відміну від C, де логічна умова може оперувати цілими числами (де \texttt{0} — це \texttt{false}, 
    а будь-яке інше число — \texttt{true}), у Java тип \texttt{boolean} є ізольованим. Конструкції 
    розгалуження приймають виключно значення \texttt{true} або \texttt{false}, що унеможливлює помилки 
    вигляду \texttt{if (a = b)}.
\end{itemize}

\subsection{Ключові відмінності в роботі з пам'яттю та об'єктами}

В межах виконання лабораторної роботи було досліджено три фундаментальні відмінності Java від мов C/C++, 
які критично впливають на проектування програм:

\begin{enumerate}
    \item \textbf{Повна відсутність вказівників:} 
    У Java повністю вилучено арифметику вказівників, прямий доступ до адрес пам'яті та 
    оператори \texttt{*}, \texttt{\&}, \texttt{->}. Усі об'єкти та масиви є посилальними 
    типами й керуються через безпечні посилання.
    \item \textbf{Автоматичне керування пам'яттю (Garbage Collection):} 
    Розробнику більше не потрібно вручну звільняти пам'ять за допомогою \texttt{free()} 
    чи \texttt{delete}. Об'єкти створюються у кучі (\texttt{heap}) за допомогою ключового 
    слова \texttt{new}, а збірник сміття (GC) автоматично утилізує їх, коли на них не залишається 
    активних посилань.
    \item \textbf{Масиви як повноцінні об'єкти:} 
    На відміну від C/C++, де масив є просто вказівником на перший елемент і вимагає окремої
     передачі його розміру у функції, в Java масив є об'єктом. Він інкапсулює у собі 
     властивість \texttt{length}, що дозволяє безпечно контролювати межі ітерацій і унеможливлює 
     помилку переповнення буфера (\texttt{ArrayIndexOutOfBoundsException}).
\end{enumerate}
